\section*{引言与实验约定}

深度模型在源域上获得较高分类准确率，并不意味着其特征能够直接迁移到成像条件不同的目标域。Deep CORAL 将源域分类损失与跨域协方差距离联合优化，为这一问题提供了结构简洁的训练目标\cite{sun2016deep}。这一设计也留下了直接的探索空间：把约束加到更多层、根据训练状态自动调整强度，或者使用更强的特征网络，是否会进一步改善迁移？当这些扩展没有形成预期收益时，问题又出在优化强度、特征表示，还是对齐的统计对象？

本文沿这条问题链组织四部分实验。第一部分比较不同层及层组合，检验增加约束位置的收益；第二部分尝试梯度范数平衡，并用固定权重日志分析训练何时应当结束；第三部分更换为 ResNet-50，观察强预训练与共享 BatchNorm 条件下 CORAL 的额外作用；第四部分转向类别条件对齐，比较自定义类条件 CORAL 与 DSAN 的 LMMD\cite{zhu2021dsan}，通过结果与数学性质反思统计目标的选择。各部分按照“探索与实验—阶段结论与反思”展开，后续探索由前一部分留下的问题引出。

全部实验使用 Office-31 的 Amazon$\to$Webcam 方向，31 类；源域包含 2,817 张图像，目标域包含 795 张图像，随机种子均为 42。目标准确率对全部 795 张图像计算，差值以百分点（pp）表示。源域准确率来自训练图片上的全量评估。本文整理截至 2026 年 10 月 12 日的 27 个归档运行，共 1,593 轮、22,968 次参数更新；本次论文整理没有新增训练。原汇报中的“Deep CORAL 原实验复现”在此指本地实现的基线实验，不表示与原论文全部训练设置逐项相同。

各阶段采用各自的结果选择方式：AlexNet 位置实验与子域对照以固定第 50 轮为主；100 轮权重实验同时呈现最终值与历史最佳值；ResNet CORAL 报告损失规则决定的停止轮，并与无 CORAL 对照的相同轮次比较；ResNet LMMD 按目标域准确率选择最佳模型。表中历史最佳值与固定轮结果分别列示。

\textbf{目标标签的用途随实验协议披露。}所有适应损失均不读取目标真实标签。AlexNet 的历史探索已查看目标评估曲线，损失早停规则的部分设计也参考了既有曲线；子域六组在阶段内固定协议与第 50 轮结果，目标真值用于评估及事后诊断；ResNet LMMD 则明确用目标真值进行最佳轮选择及 30 轮无提升停止。因此，后者的 93.71\% 属于目标标签参与模型选择的结果。所有结果均为单种子、单迁移方向的观察，本文据此比较具体配置和提出机制解释，不估计跨种子显著性。
